% \iffalse
%<*driver>
\ProvidesFile{seqenseigne-flashcard.dtx}%
[2025/03/16 v1.0-dev seqenseigne – Module flashcard (cartes mémoire imprimables)]
%</driver>
% \fi
%
% =============================================================================
% seqenseigne-flashcard.dtx
% Module : FLASHCARD
% Rôle   :
%   - Définir le format de page pour les cartes mémoire
%   - Fournir des environnements pour créer des cartes recto/verso
%   - Générer automatiquement des séries de cartes depuis les bases de données
%
% Dépendances :
%   - seqenseigne-theme.dtx  (couleurs, boîtes tcolorbox)
%   - seqenseigne-data.dtx   (accès aux objectifs et connaissances)
%   - paquet 'pgf' ou 'tikz' (déjà chargé par core)
%
% Formats de page disponibles (option du paquet) :
%   portrait  : A6 portrait (105×148 mm), 4 cartes par feuille A4
%   paysage   : demi-A4 paysage (148×105 mm), 2 cartes par feuille A4
%
% Structure d'une carte :
%   Recto : bandeau thème/séquence, objectif centré, zone nom
%   Verso : trois blocs de critères (TB/S/F) + zone exemples libre
%
% Usage sans données CSV (mode manuel) :
%   \begin{seqFlashcard}[options]{contenu verso}%
%     contenu recto (optionnel, sinon généré)
%   \end{seqFlashcard}%
%
% Usage avec données CSV (mode automatique) :
%   \seqFlashcardsSequence{codeNiveau}{codeSequence}%
%   \seqFlashcardsAnnee{codeNiveau}%
%
% IMPORTANT — Ce module ne peut PAS être chargé dans le même document
% qu'un livret ou une évaluation : il redéfinit la géométrie de page.
% Créer un fichier .tex séparé dédié aux flashcards.
% =============================================================================
%
%<*package>
\NeedsTeXFormat{LaTeX2e}%
\ProvidesPackage{seqenseigne-flashcard}%
[2025/03/16 v1.0-dev seqenseigne flashcard module]

% =============================================================================
% Dépendances
% =============================================================================
% tcolorbox et ses bibliothèques sont déjà chargés par seqenseigne-theme.
% On ne charge que ce qui est propre à ce module.
\RequirePackage{xkeyval}%
\RequirePackage{xifthen}%
\RequirePackage{etoolbox}%
\RequirePackage{geometry}%
\RequirePackage{everypage}  % pour alterner recto/verso sans erreur pdflatex

% =============================================================================
% Options du module
% =============================================================================
%
% \begin{macro}{format}%
% Option 'format' : portrait (A6, 4 cartes/feuille) ou paysage (demi-A4, 2/feuille).
% Défaut : portrait.
% \begin{macrocode}%
\newif\ifseqfc@paysage\seqfc@paysagefalse
\DeclareOption{portrait}{\seqfc@paysagefalse}%
\DeclareOption{paysage}{\seqfc@paysagetrue}%
\DeclareOption{ltxdoc}{}%
\ProcessOptions\relax
% \end{macrocode}%
% \end{macro}%

% =============================================================================
% Géométrie de page
% =============================================================================
%
% \begin{macro}{\seqfc@setupgeometry}%
% Configure la géométrie selon le format choisi.
% \begin{macrocode}%
\newcommand{\seqfc@setupgeometry}{%
    \ifseqfc@paysage
        % Demi-A4 paysage : 148 × 105 mm — 2 cartes par feuille A4
        \newgeometry{%
            paperwidth  = 148mm,
            paperheight = 105mm,
            left        = 6mm,
            right       = 6mm,
            top         = 6mm,
            bottom      = 6mm,
            nohead,
            nofoot
        }%
    \else
        % A6 portrait : 105 × 148 mm — 4 cartes par feuille A4
        \newgeometry{%
            paperwidth  = 105mm,
            paperheight = 148mm,
            left        = 5mm,
            right       = 5mm,
            top         = 5mm,
            bottom      = 5mm,
            nohead,
            nofoot
        }%
    \fi
}%
% \end{macrocode}%
% \end{macro}%

% Appliquer la géométrie dès le chargement du module
\AtBeginDocument{\seqfc@setupgeometry}%

% =============================================================================
% Compteur de cartes
% =============================================================================
\newcounter{seqfc@num}%
\setcounter{seqfc@num}{0}%

% Macro interne : numéro courant formaté
\newcommand{\seqfc@numcourant}{\arabic{seqfc@num}}%

% =============================================================================
% Couleurs et styles internes
% =============================================================================
%
% Les couleurs TB/S/F sont celles utilisées dans tout le paquet pour la maîtrise.
% On les redéfinit localement pour éviter toute dépendance sur les variables
% de thème (\seq@couleuri@vif etc.) qui peuvent ne pas être initialisées
% si seqenseigne-theme n'est pas chargé.
%
% Si seqenseigne-theme EST chargé (usage normal), ces définitions sont remplacées
% par les couleurs du thème courant via \seqfc@applytheme.
%
\providecommand{\seqfc@couleur@tb}{green!60!black}%
\providecommand{\seqfc@couleur@s}{green!60!white}%
\providecommand{\seqfc@couleur@f}{yellow!80!black}%
\providecommand{\seqfc@couleur@titre}{blue!70!black}%
\providecommand{\seqfc@couleur@fond}{blue!10!white}%
\providecommand{\seqfc@couleur@cadre}{blue!70!black}%

% \begin{macro}{\seqfc@applytheme}%
% Synchronise les couleurs de flashcard avec le thème seqenseigne courant.
% À appeler après \seqSetColors ou \seq@setColors.
% \begin{macrocode}%
\newcommand{\seqfc@applytheme}{%
    \renewcommand{\seqfc@couleur@titre}{\seq@couleuri@vif}%
    \renewcommand{\seqfc@couleur@fond}{\seq@couleurii@pale}%
    \renewcommand{\seqfc@couleur@cadre}{\seq@couleuri@vif}%
}%
% \end{macrocode}%
% \end{macro}%

% =============================================================================
% Macros de mise en page interne
% =============================================================================

% \begin{macro}{\seqfc@recto}%
% Compose le recto d'une carte.
% #1 : code niveau
% #2 : code séquence
% #3 : code objectif
% #4 : libellé de l'objectif (texte)
% #5 : libellé de la séquence (texte)
% \begin{macrocode}%
\newcommand{\seqfc@recto}[5]{%
    \stepcounter{seqfc@num}%
    \noindent
    \begin{tcolorbox}[
        height          = \textheight,
        colback         = \seqfc@couleur@fond,
        colframe        = \seqfc@couleur@cadre,
        colbacktitle    = \seqfc@couleur@titre,
        coltitle        = white,
        fonttitle       = \bfseries\small,
        title           = {%
            \seqNiveauGetNomCourt{#1}\hfill
            #2~--~#5\hfill
            Fiche n°\,\seqfc@numcourant
        },
        enhanced,
        attach boxed title to top left={yshift=-4pt, xshift=10pt},
        top             = 14pt,
        bottom          = 4pt,
    ]
        % Zone objectif : centré verticalement, texte en grand
        \vfill
        \begin{center}%
            {\Large\bfseries\color{\seqfc@couleur@cadre} Objectif~#3}\par
            \vspace{6pt}%
            {\large #4}%
        \end{center}%
        \vfill
        % Zone nom élève en bas
        \tcblower
        \small
        \textit{Version de \datemoisannee\today}\hfill
        Nom~: \underline{\hspace{5cm}}%
    \end{tcolorbox}%
    \newpage
}%
% \end{macrocode}%
% \end{macro}%

% \begin{macro}{\seqfc@verso}%
% Compose le verso d'une carte avec les critères de maîtrise et une zone libre.
% #1 : critère Très Bon (MaitriseTB)
% #2 : critère Satisfaisant (MaitriseS)
% #3 : critère À consolider (MaitriseF)
% #4 : contenu libre (exemples, formules, schémas — peut être vide)
% \begin{macrocode}%
\newcommand{\seqfc@verso}[4]{%
    \noindent
    \begin{tcolorbox}[
        height          = \textheight,
        colback         = white,
        colframe        = \seqfc@couleur@cadre,
        colbacktitle    = \seqfc@couleur@titre,
        coltitle        = white,
        fonttitle       = \bfseries\small,
        title           = {Critères de maîtrise — Fiche n°\,\seqfc@numcourant},
        enhanced,
        attach boxed title to top left={yshift=-4pt, xshift=10pt},
        top             = 14pt,
        bottom          = 4pt,
    ]
        % Bloc Très Bon
        \begin{tcolorbox}[
            colback  = \seqfc@couleur@tb,
            colframe = \seqfc@couleur@tb,
            coltitle = white,
            fonttitle= \bfseries\footnotesize,
            title    = {Très bon},
            top=2pt, bottom=2pt, left=4pt, right=4pt,
        ]
            \footnotesize #1
        \end{tcolorbox}%
        \vspace{3pt}%
        % Bloc Satisfaisant
        \begin{tcolorbox}[
            colback  = \seqfc@couleur@s,
            colframe = \seqfc@couleur@s,
            coltitle = black,
            fonttitle= \bfseries\footnotesize,
            title    = {Satisfaisant},
            top=2pt, bottom=2pt, left=4pt, right=4pt,
        ]
            \footnotesize #2
        \end{tcolorbox}%
        \vspace{3pt}%
        % Bloc À consolider
        \begin{tcolorbox}[
            colback  = \seqfc@couleur@f,
            colframe = \seqfc@couleur@f,
            coltitle = black,
            fonttitle= \bfseries\footnotesize,
            title    = {\`A consolider},
            top=2pt, bottom=2pt, left=4pt, right=4pt,
        ]
            \footnotesize #3
        \end{tcolorbox}%
        % Zone exemples libre (si fournie)
        \ifthenelse{\isempty{#4}}{}{%
            \vspace{4pt}%
            \begin{tcolorbox}[
                colback  = white,
                colframe = \seqfc@couleur@cadre,
                coltitle = white,
                colbacktitle = \seqfc@couleur@cadre,
                fonttitle= \bfseries\footnotesize,
                title    = {Exemples},
                top=2pt, bottom=2pt, left=4pt, right=4pt,
                enhanced,
                attach boxed title to top left={yshift=-3pt, xshift=8pt},
            ]
                \footnotesize #4
            \end{tcolorbox}%
        }%
        \vfill
    \end{tcolorbox}%
    \newpage
}%
% \end{macrocode}%
% \end{macro}%

% =============================================================================
% Interface publique
% =============================================================================

% ----------------------------------------------------------------------------
% \begin{macro}{\seqFlashcard}%
% \IndexMacro{\seqFlashcard}%
% Crée une paire recto/verso pour UN objectif spécifié manuellement.
% Mode sans données CSV : tous les contenus sont fournis en arguments.
%
% Utilisation :
%   \seqFlashcard
%     [niveau=N11, sequence=S03, objectif=02]
%     {libellé objectif}%
%     {libellé séquence}%
%     {critère TB}%
%     {critère S}%
%     {critère F}%
%     {exemples (peut être vide {})}%
%
% #1 (optionnel) : options clé=valeur (niveau, sequence, objectif)
% #2 : libellé de l'objectif (recto)
% #3 : libellé de la séquence (recto, bandeau)
% #4 : critère Très Bon (verso)
% #5 : critère Satisfaisant (verso)
% #6 : critère À consolider (verso)
% #7 : exemples/contenu libre du verso (peut être {})
% \begin{macrocode}%
\define@cmdkey[seqfc]{flashcard}{objectif}[00]{}%
\newcommand{\seqFlashcard}[7][]{%
    \setkeys[seqfc]{flashcard}{niveau,sequence,objectif,#1}%
    \seqfc@recto
        {\cmdseqfc@flashcard@niveau}%
        {\cmdseqfc@flashcard@sequence}%
        {\cmdseqfc@flashcard@objectif}%
        {#2}%
        {#3}%
    \seqfc@verso{#4}{#5}{#6}{#7}%
}%
% \end{macrocode}%
% \end{macro}%

% ----------------------------------------------------------------------------
% \begin{macro}{\seqFlashcardAuto}%
% \IndexMacro{\seqFlashcardAuto}%
% Crée une paire recto/verso pour UN objectif, en lisant les données des bases.
% Requiert \seqLoadData (ou \seqSetDataPath + \seqLoadData) préalablement appelé.
%
% #1 (optionnel) : options clé=valeur
%   niveau   : code niveau   (ex. N11)
%   sequence : code séquence (ex. S03)
%   objectif : code objectif (ex. 02)
%   exemples : contenu libre du verso (texte LaTeX, peut être vide)
%
% Usage :
%   \seqFlashcardAuto[niveau=N11, sequence=S03, objectif=02, exemples={}]
% \begin{macrocode}%
\define@cmdkey[seqfc]{flashcardauto}{objectif}[00]{}%
\define@cmdkey[seqfc]{flashcardauto}{exemples}[]{}%
\newcommand{\seqFlashcardAuto}[1][]{%
    \setkeys[seqfc]{flashcardauto}{niveau,sequence,objectif,exemples,#1}%
    % Synchroniser le thème avec la séquence courante
    \seqSetColors%
    \seqfc@applytheme
    % Recto
    \seqfc@recto
        {\cmdseqfc@flashcardauto@niveau}%
        {\cmdseqfc@flashcardauto@sequence}%
        {\cmdseqfc@flashcardauto@objectif}%
        {\seqObjectifGetNom
            {\cmdseqfc@flashcardauto@niveau}%
            {\cmdseqfc@flashcardauto@sequence}%
            {\cmdseqfc@flashcardauto@objectif}}%
        {\seqSequenceGetNom
            {\cmdseqfc@flashcardauto@niveau}%
            {\cmdseqfc@flashcardauto@sequence}}%
    % Verso
    \seqfc@verso
        {\seqObjectifGetMaitriseTB
            {\cmdseqfc@flashcardauto@niveau}%
            {\cmdseqfc@flashcardauto@sequence}%
            {\cmdseqfc@flashcardauto@objectif}}%
        {\seqObjectifGetMaitriseS
            {\cmdseqfc@flashcardauto@niveau}%
            {\cmdseqfc@flashcardauto@sequence}%
            {\cmdseqfc@flashcardauto@objectif}}%
        {\seqObjectifGetMaitriseF
            {\cmdseqfc@flashcardauto@niveau}%
            {\cmdseqfc@flashcardauto@sequence}%
            {\cmdseqfc@flashcardauto@objectif}}%
        {\cmdseqfc@flashcardauto@exemples}%
}%
% \end{macrocode}%
% \end{macro}%

% ----------------------------------------------------------------------------
% \begin{macro}{\seqFlashcardsSequence}%
% \IndexMacro{\seqFlashcardsSequence}%
% Génère automatiquement TOUTES les cartes des objectifs d'une séquence.
% Boucle sur seq@objectifs@sequence après filtrage.
% Requiert \seqLoadData préalablement appelé.
%
% #1 : code niveau   (ex. N11)
% #2 : code séquence (ex. S03)
%
% Usage :
%   \seqFlashcardsSequence{N11}{S03}%
% \begin{macrocode}%
\newcommand{\seqFlashcardsSequence}[2]{%
    % Fixer le contexte niveau/séquence
    \seqSetCodeNiveau{#1}%
    \seqSetCodeSequence{#2}%
    % Filtrer les données et synchroniser le thème
    \seq@bases@filter{\seq@code@niveau}{\seq@code@sequence}%
    \seqSetColors%
    \seqfc@applytheme
    % Boucler sur tous les objectifs de la séquence
    \DTLforeach*{seq@objectifs@sequence}{%
        \seqfc@tmpCode=CodeObjectif,%
        \seqfc@tmpNom=Nom,%
        \seqfc@tmpTB=MaitriseTB,%
        \seqfc@tmpS=MaitriseS,%
        \seqfc@tmpF=MaitriseF%
    }{%
        \seqfc@recto
            {\seq@code@niveau}%
            {\seq@code@sequence}%
            {\seqfc@tmpCode}%
            {\seqfc@tmpNom}%
            {\seqSequenceGetNom}%
        \seqfc@verso
            {\seqfc@tmpTB}%
            {\seqfc@tmpS}%
            {\seqfc@tmpF}%
            {}%  verso sans exemples en génération automatique
    }%
}%
% \end{macrocode}%
% \end{macro}%

% ----------------------------------------------------------------------------
% \begin{macro}{\seqFlashcardsAnnee}%
% \IndexMacro{\seqFlashcardsAnnee}%
% Génère toutes les cartes de toutes les séquences d'un niveau (toute l'année).
% Boucle sur la base des séquences du cycle correspondant.
% Requiert \seqLoadData préalablement appelé.
%
% #1 : code niveau (ex. N11)
%
% Usage :
%   \seqFlashcardsAnnee{N11}%
% \begin{macrocode}%
\newcommand{\seqFlashcardsAnnee}[1]{%
    % Fixer le niveau courant
    \seqSetCodeNiveau{#1}%
    % Récupérer le code cycle du niveau
    \seqDTLfetchsave{\seqfc@tmp@cycle}{param_niveaux}{Code}{\seq@code@niveau}{CodeCycle}%
    % Boucler sur les séquences du cycle
    \DTLforeach*{\seqfc@tmp@cycle _sequences}{%
        \seqfc@tmpSeq=Code%
    }{%
        \seqFlashcardsSequence{\seq@code@niveau}{\seqfc@tmpSeq}%
    }%
}%
% \end{macrocode}%
% \end{macro}%

% ----------------------------------------------------------------------------
% \begin{macro}{\seqFlashcardVersoBlanc}%
% \IndexMacro{\seqFlashcardVersoBlanc}%
% Insère un verso blanc (pour impression recto-verso sur du papier ordinaire
% quand la carte précédente avait un verso, mais la suivante commence un recto).
% Utile pour forcer l'alignement recto/verso lors d'une impression deux faces.
% \begin{macrocode}%
\newcommand{\seqFlashcardVersoBlanc}{%
    \noindent
    \begin{tcolorbox}[
        height   = \textheight,
        colback  = white,
        colframe = white,
    ]
    \end{tcolorbox}%
    \newpage
}%
% \end{macrocode}%
% \end{macro}%

% ----------------------------------------------------------------------------
% \begin{macro}{\seqFlashcardResetCompteur}%
% \IndexMacro{\seqFlashcardResetCompteur}%
% Remet le compteur de cartes à zéro.
% Utile pour recommencer la numérotation entre deux séries.
% \begin{macrocode}%
\newcommand{\seqFlashcardResetCompteur}{%
    \setcounter{seqfc@num}{0}%
}%
% \end{macrocode}%
% \end{macro}%

% =============================================================================
% Variables temporaires internes (déclarées globalement pour éviter
% les problèmes de portée dans les boucles DTL)
% =============================================================================
\newcommand{\seqfc@tmpCode}{}%
\newcommand{\seqfc@tmpNom}{}%
\newcommand{\seqfc@tmpTB}{}%
\newcommand{\seqfc@tmpS}{}%
\newcommand{\seqfc@tmpF}{}%
\newcommand{\seqfc@tmp@cycle}{}%

% Fin du module flashcard
%</package>
